% ACM conference paper (sigconf), on the acmart class that comes with TeX Live.
% Guidelines: https://www.acm.org/publications/proceedings-template
% For review, use \documentclass[sigconf,review,anonymous]{acmart}.
\documentclass[sigconf]{acmart}

\setcopyright{acmlicensed}
\copyrightyear{2026}
\acmYear{2026}
\acmDOI{XXXXXXX.XXXXXXX}
\acmConference[Conference acronym 'XX]{Make sure to enter the correct conference title from your rights confirmation email}{June 03--05, 2026}{Woodstock, NY}
\acmISBN{978-1-4503-XXXX-X/2026/06}

\begin{document}

\title{Paper Title}

\author{First Author}
\affiliation{%
  \institution{Organization}
  \city{City}
  \country{Country}}
\email{email@example.com}

\author{Second Author}
\affiliation{%
  \institution{Organization}
  \city{City}
  \country{Country}}
\email{email@example.com}

\begin{abstract}
A short summary of the problem, the approach and the results, in one paragraph.
\end{abstract}

\begin{CCSXML}
<ccs2012>
 <concept>
  <concept_id>10010147.10010257</concept_id>
  <concept_desc>Computing methodologies~Machine learning</concept_desc>
  <concept_significance>500</concept_significance>
 </concept>
</ccs2012>
\end{CCSXML}
\ccsdesc[500]{Computing methodologies~Machine learning}

\keywords{first keyword, second keyword, third keyword}

\maketitle

\section{Introduction}
Say what the problem is and why it matters, then what this paper does~\cite{lamport1994}.

\section{Method}
\begin{equation}
  y = f(x) + \epsilon
  \label{eq:model}
\end{equation}

\begin{figure}[h]
  \centering
  \includegraphics[width=\linewidth]{example-image}
  \caption{A figure that spans one column.}
  \Description{A placeholder image.}
  \label{fig:example}
\end{figure}

\section{Experiments}
Results go here; see Figure~\ref{fig:example}.

\section{Conclusion}
What was shown, and what comes next.

\begin{acks}
Funding and people to thank.
\end{acks}

\bibliographystyle{ACM-Reference-Format}
\bibliography{references}

\end{document}
